% Nouveau EGA V (cote 25): archive cote 25, site page 112 (PDF page 113).
% typescript. Draft compared: none.
% AI-transcribed by Claude Opus 5.5 (claude-opus-5-5, adaptive thinking, effort high), Message Batches API, batch ega25b (msgbatch_01A7n5YuXubUPNK1zsepi8RG), 2026-11-13; not yet checked by eye.
\scanpage{25/112}
\mspage{15}

\deleted{\illegible{1}} \margin{\fbox{N$^{\text{os}}$ 3,4}} \hand{3)} \inserted{\uline{Systèmes linéaires de diviseurs, et \unsure{morphismes} dans les fibrés \unsure{projectifs}.}}
\deleted{16. Traductions en langage de systèmes linéaires.}

\deleted{Soit X plat et loc de présentation finie sur S, et T un S-préschéma.}
\deleted{On appelle \uline{famille de} \xout{xxxxxxx} \uline{diviseurs} (sous-entendu: de Cartier)}
\deleted{\margin{\unsure{r}}\uline{positifs}\xout{x}\hand{\unsure{$^{3}$}} \inserted[typed]{\uline{de X/S} \hand{,}} \hand{/} \uline{paramétrée par} S, un sous-préschéma fermé de $X\times_S T$ qui est}
\deleted{\inserted{un diviseur de Cartier rel. à T, i.e. qui est} \hand{(}intersection complète \xout{xxxxxxxxxxx} relative de codimension 1 relativement}
\deleted{à T en chacun de ses points, ou \xout{xxx} ce qui revient au même, qui est plat}
\deleted{\inserted[typed]{loc de prés finie} \hand{(}par rapport à T, et induit sur chaque fibre de $X_{T\,t} = X_{s\,k(t)}$ de $X_T$ sur}
\deleted{T un ``diviseur de Cartier positif'' i.e. un sous-préschéma qui est en}
\deleted{chaque \inserted{de ses} point intersection complète \inserted[typed]{dans $X_t$} de codimension 1. L'ensemble}
\deleted{$\mathrm{Div}_{X/S}(T)$ de telles familles de \inserted[typed]{di}viseurs est, pour T variable, un}
\deleted{foncteur contravar\xout{i}iant à valeurs dans (Ens), noté $\mathrm{Div}_{X/S}$ . Nous verrons}
\deleted{au Chap V qu'il est représentable lorsque $f:X\quad S$ est projectif. Lorsque}
\deleted{X = P(E) fibré projectif défini par E localement libre de type fini,}
\deleted{il est possible avec les seuls résultats du Chap III de prouver que}
\deleted{$\mathrm{Div}_{P/S}$ est représentable par le préschéma somme des $P(\mathrm{Symm}^n(\underline{E}))$ , les}
\deleted{composants de cette somme correspondant aux diviseurs de degré relatif}
\deleted{n , \inserted{\unsure{\S\ n$^{\circ}$2}} Ce résultat n'est pas indispensable pour la suite, et il suffira de}
\deleted{le signaler en remarque; nous l'inclurons peut-être quelque part dans}
\deleted{III, et au plus tard dans V .}

Soit D une famille de diviseurs \inserted[typed]{sur X/S} paramétrée par T (nous sous-enten%
% NOTE: the end of this line is cut off at the right edge of the scan
dons par la suite ``positif''). Un point $x\in X$ est dit ``point fixe'' \xout{pour}pour
cette famille de diviseurs si $pr_1^{-1}(x)\subset D$ \inserted{ensemblistement}, de sorte que l'ensemble des
points non fixes est le complémentaire de $pr_1(\xout{X}X\times_S T - D)$; par suite, si
$T\longrightarrow S$ est universellement ouvert (par exemple plat loc de prés finie)
l'ensemble des points fixes est fermé. On dit que la famille de diviseurs
est ``sans points fixes'' si l'ensemble \inserted{$Z$} des points fixes est vide.

\mspage{112}
